\section{Further general obstructions}
\label{sec:theoretical-oct16}
The following arguments exclude nine further source tensors without a
pivotality assumption. The proofs use a conductor bound, a
Frobenius--Perron index, cyclic-module divisibility, and changes of Morita representative.

\begin{proposition}\label{prop:r8-conductor-four}
The rings $R8\_2\_35$, $R8\_2\_38$, $R8\_2\_346$ and
$R8\_2\_406$ have no categorification over $\C$.
\end{proposition}
\begin{proof}
Each tensor has four invertibles and four simples of FP dimension
$d=3+\sqrt{10}$. By the published \cite[Lemma~3.1]{SchNear}, existence
of any categorification implies existence of a pseudo-unitary Galois
conjugate, whose canonical spherical structure has global dimension
\[
 D=4(1+d^2)=80+24\sqrt{10},\qquad DD'=640.
\]
Its simple-dimension and global-dimension fields both equal
$\Q(\sqrt{10})$. The nontrivial quadratic character has conductor $40$,
so \cite[Theorem~5.1]{SchDimensionBlock} requires
\[
 \left|\frac1D-\frac1{D'}\right|^2\leq\frac1{40}.
\]
The left side is $9/160>1/40$, a contradiction.
\end{proof}

\begin{proposition}\label{prop:r8-209-index}
The ring $R8\_2\_209$ has no categorification over $\C$.
\end{proposition}
\begin{proof}
Put $d=(5+\sqrt{29})/2$. The source dimension vector is
$(1,1,1,d,d-1,d+1,d,d)$, and its first three basis elements form a
pointed subcategory of dimension $3$. Hence
\[
 \FPdim(\mathcal C)=5(1+d^2)=\frac{145+25\sqrt{29}}2.
\]
By \cite[Proposition~8.15]{ENO}, its quotient by $3$ must be an
algebraic integer. Its field trace is $145/3$, which is not an integer.
\end{proof}

% Standalone common lemma; copy/input once per rank manuscript.
% Requires EGODimPQ bibliography entry and standard amsthm/math notation.
\begin{lemma}[Cyclic-module divisibility]\label{lem:cyclic-module-divisibility}
Let $H\cong C_p$, with $p$ prime, be a group of invertibles in a
fusion category over $\mathbb C$. Suppose that simple objects $X,Y,Z$
satisfy $hX\cong X$, $Yh\cong Y$, and $hZ\cong Z\cong Zh$ for every
$h\in H$. If $\operatorname{FPdim}(Z)/p$ is not an algebraic integer,
then $p\mid N_{XY}^{Z}$.
\end{lemma}
\begin{proof}
The rank-one module category generated by $X$ trivializes the pointed
associator on $H$. Thus $A=\bigoplus_{h\in H}h$ admits its regular
connected separable algebra structure. Since
$H^2(C_p,\mathbb C^\times)=0$, the stabilizer cocycles can be
trivialized to give $X$ a left $A$-module structure and $Y$ a right
$A$-module structure. Hence $X\otimes Y$ is an $A$-bimodule.

The $A$-bimodules supported on the fixed simple $Z$ correspond to
projective representations of $C_p^2$: the projective-representation
dimension equals the underlying multiplicity of $Z$. A scalar
two-cocycle on $C_p^2$ has either zero or nondegenerate commutator
pairing. Its twisted group algebra is accordingly
$\mathbb C^{p^2}$ or $M_p(\mathbb C)$. The first possibility would
give a simple bimodule of FP dimension
$\operatorname{FPdim}(Z)/\operatorname{FPdim}(A)
=\operatorname{FPdim}(Z)/p$, contrary to algebraic integrality.
In the second possibility every simple bimodule supported on $Z$
has underlying object $pZ$. Decomposing the $Z$-isotypic bimodule
summand of $X\otimes Y$ proves the assertion. This is the regular
cyclic-algebra method of \cite[Lemmas~6.1--6.2]{EGODimPQ}, with
the stabilizer cocycle retained rather than assumed trivial.
\end{proof}

For later use, a free left $C_p$ orbit whose right orbit agrees with
it has cyclic left-right stabilizer of order $p$. Its block in
$A$-\emph{bimod}$(\mathcal C)$ therefore consists of $p$ simples
of the original FP dimension. A label fixed on both sides instead
contributes either $p^2$ simples of dimension $\FPdim(X)/p$ or
one simple of dimension $\FPdim(X)$, by the preceding proof.

\begin{proposition}\label{prop:r8-260-morita}
The ring $R8\_2\_260$ has no categorification over $\C$.
\end{proposition}
\begin{proof}
Put $d=3+2\sqrt3$. Its dimension vector is
$(1,1,1,d,d,d,d,3)$, and the first three basis elements form $C_3$.
They fix label $3$ on both sides. Since
$\operatorname{Norm}(d/3)=-1/3$, Lemma~\ref{lem:cyclic-module-divisibility}
requires $3\mid N_{33}^{3}$. The source tensor has $N_{33}^{3}=2$.
\end{proof}

\begin{proposition}\label{prop:r8-380-673-morita}
The rings $R8\_2\_380$ and $R8\_2\_673$ have no categorification
 over $\C$.
\end{proposition}
\begin{proof}
Write $x=4+\sqrt6$ and $y=3+2\sqrt6$. Each ring has a subcategory
$\mathcal E$ with the representation ring of the dihedral group
$D_{10}$ of order ten. It occupies labels $(0,3,4,5)$ in $R8\_2\_380$
and $(0,1,6,7)$ in $R8\_2\_673$. Its nontrivial invertible $g$ fixes
the two dimension-two labels, so the regular algebra $A=1+g$ exists.
By \cite[Lemma~6.2 and the proof of Theorem~6.3]{EGODimPQ},
$A$-\emph{bimod}$(\mathcal E)$ is pointed with group $D_{10}$.

The outside labels form two free $g$-orbits, of dimensions $x$ and
$y$, with the same left and right actions. They are $(1,2),(6,7)$
in the first tensor and $(2,5),(3,4)$ in the second. The cyclic
stabilizer calculation gives two $x$-dimensional and two
$y$-dimensional simples in the dual category
$\mathcal B=A$-\emph{bimod}$(\mathcal C)$; all its other simples
are the ten invertibles of $D_{10}$.

The rotation subgroup $H=C_5$ fixes every noninvertible on both
sides, since each dimension fiber has only two elements. Moreover
\[
 \operatorname{Norm}(x/5)=\frac25,\qquad
 \operatorname{Norm}(y/5)=-\frac35.
\]
For an $x$-dimensional simple $X$, Lemma~\ref{lem:cyclic-module-divisibility}
therefore makes every noninvertible multiplicity in $XX^*$ divisible
by $5$. Write $x^2=a+bx+cy$ for the resulting dimension equation,
where $a,b,c$ are nonnegative integers. The $\sqrt6$ coefficient gives
$b+2c=8$, contradicting $5\mid b,c$.
\end{proof}

\begin{proposition}\label{prop:r8-745-morita}
The ring $R8\_2\_745$ has no categorification over $\C$.
\end{proposition}
\begin{proof}
Put $d=1+\sqrt2$ and let $g$ be label $1$. The dimension vector is
$(1,1,d,d,2,2,2d,2d)$. The subsets
$\mathcal D=(0,1,2,3)$ and $\mathcal E=(0,1,4,5)$ are fusion
subcategories, the latter with the representation ring of the dihedral
group $D_{10}$ of order $10$. The element $g$ exchanges labels $2,3$
and fixes labels $4,5,6,7$ on both sides. Thus $A=1+g$ exists.

In $\mathcal B=A$-\emph{bimod}$(\mathcal C)$, the subcategory
$\mathcal E'=A$-\emph{bimod}$(\mathcal E)$ is pointed with group
$D_{10}$, by \cite[Lemma~6.2 and the proof of Theorem~6.3]{EGODimPQ}.
The subcategory $\mathcal D'=A$-\emph{bimod}$(\mathcal D)$ has two
invertibles and two simples of dimension $d$, by the cyclic-stabilizer
calculation above. Each of labels $6,7$ contributes either four
simples of dimension $d$ or one of dimension $2d$. These are all the
simples of $\mathcal B$, so it has precisely ten invertibles and
$2$, $6$ or $10$ simples of dimension $d$.

For a $d$-dimensional simple $X$, the dimension of $XX^*$ is
$a+bd$ with $a,b\in\NN$, where $a$ is its invertible stabilizer order.
The identity $d^2=1+2d$ forces $a=1$. Thus the $D_{10}$ action on
these simples is free, and their number must be $10$. It follows that
$\mathcal B$ is generalized near-group, with no $2d$-dimensional
simples remaining.

By \cite[Proposition~3.6(1)]{DongSunNear}, the adjoint
subcategory $\mathcal B_{\mathrm{ad}}$ is nonpointed and is contained
in every nonpointed fusion subcategory, hence in $\mathcal D'$.
There is no proper nonpointed fusion subcategory of $\mathcal D'$:
one noninvertible would give the forbidden rank-two rule
$X^2=1+2X$ \cite{Ostrik2}, while two noninvertibles force the other
invertible by reciprocity. Therefore
$\mathcal B_{\mathrm{ad}}=\mathcal D'$. Its order-two invertible group
must be normal in $D_{10}$, since conjugation preserves the adjoint
subcategory. No such normal subgroup exists.
\end{proof}

The source-bound verifiers in
\sha{evidence/theoretical_20261016/rank8} check all dimension
characters and subcategory tensors used above, including the indicated
left and right actions. The categorical implications are the preceding
proofs and their explicitly cited dependencies.
